% 中文 SCI 论文 LaTeX 模板（通用学术格式）
% 适用：多数 SCI 期刊投稿前的预印本 / 内部版本
% 编译方式：xelatex -> bibtex -> xelatex*2
\documentclass[12pt,a4paper]{article}

% ========== 基础宏包 ==========
\usepackage{ctex}                       % 中文支持
\usepackage[margin=2.5cm]{geometry}     % 页面边距
\usepackage{graphicx}                   % 插图
\usepackage{booktabs}                   % 三线表
\usepackage{multirow}                   % 表格合并
\usepackage{amsmath,amssymb,amsfonts}   % 数学公式
\usepackage{algorithm}                  % 算法环境
\usepackage{algpseudocode}              % 算法伪代码
\usepackage{subcaption}                 % 子图
\usepackage{xcolor}                     % 颜色
\usepackage{hyperref}                   % 超链接
\usepackage{lineno}                     % 行号（审稿时常用）
\linenumbers

% ========== 参考文献设置 ==========
\usepackage[numbers,sort&compress]{natbib}
\bibliographystyle{unsrtnat}  % 可替换为期刊提供的 bst 文件

% ========== 标题与作者信息 ==========
\title{\textbf{论文标题：请在此处填写你的论文标题}}
\author{
    作者一$^{1,*}$, \quad
    作者二$^{2}$, \quad
    通讯作者$^{1,\dagger}$ \\[0.5em]
    \small $^{1}$ 单位一名称，城市，邮编，国家 \\
    \small $^{2}$ 单位二名称，城市，邮编，国家 \\
    \small $^{*}$ 第一作者邮箱；$^{\dagger}$ 通讯作者邮箱
}
\date{}

\begin{document}

\maketitle

% ========== 摘要 ==========
\begin{abstract}
请在此处撰写中文摘要（通常 200--300 字）。摘要应包含研究背景、目的、方法、主要结果和结论。避免使用公式、图表和参考文献编号。
\end{abstract}

\noindent\textbf{关键词：} 关键词1；关键词2；关键词3；关键词4

\vspace{1em}
\noindent\rule{\textwidth}{0.4pt}

% ========== 1. 引言 ==========
\section{引言（Introduction）}
引言部分应回答以下问题：研究背景是什么？现有工作存在哪些不足？本文的研究动机和主要贡献是什么？

例如：近年来，深度学习在计算机视觉领域取得了显著进展\cite{he2016deep}。然而，在 xxx 任务中，现有方法仍面临 xxx 问题\cite{vaswani2017attention}。针对上述问题，本文提出了一种 xxx 方法，主要贡献如下：
\begin{itemize}
    \item 贡献一：请简要描述。
    \item 贡献二：请简要描述。
    \item 贡献三：请简要描述。
\end{itemize}

% ========== 2. 相关工作 ==========
\section{相关工作（Related Work）}
按主题分类综述相关研究，突出与本工作的区别。

\subsection{子领域一}
[请填写相关研究综述。]

\subsection{子领域二}
[请填写相关研究综述。]

% ========== 3. 方法 ==========
\section{方法（Methodology）}
详细描述所提出的方法。建议包含：总体框架图、符号定义、模型结构、损失函数、训练策略等。

\subsection{问题定义}
给出任务的数学定义。例如，给定训练数据集 $\mathcal{D}=\{(x_i, y_i)\}_{i=1}^{N}$，其中 $x_i \in \mathbb{R}^{H \times W \times 3}$ 为输入图像，$y_i \in \{0, 1\}$ 为标签，目标是学习一个映射函数 $f: x \mapsto y$。

\subsection{网络结构}
请在此处插入网络结构图，并描述各模块功能：
\begin{equation}
    \mathcal{L}_{\text{total}} = \mathcal{L}_{\text{cls}} + \lambda \mathcal{L}_{\text{aux}}.
\end{equation}

\subsection{损失函数}
详细说明各项损失函数的设计动机与公式。

\subsection{训练细节}
包括优化器、学习率、批量大小、数据增强、硬件环境等。

% ========== 4. 实验 ==========
\section{实验（Experiments）}
\subsection{数据集与评价指标}
描述数据集划分、数据预处理方式和评价指标。例如：
\begin{itemize}
    \item 数据集：训练集 / 验证集 / 测试集数量。
    \item 评价指标：准确率（Accuracy）、AUC、F1-score 等。
\end{itemize}

\subsection{实现细节}
实验环境、超参数设置、复现性说明。

\subsection{与先进方法的对比}
将实验结果整理成表格。示例如表~\ref{tab:comparison}。

\begin{table}[htbp]
\centering
\caption{与现有方法的对比结果}
\label{tab:comparison}
\begin{tabular}{lcccc}
\toprule
方法 & 年份 & Acc & AUC & F1 \\
\midrule
方法A & 2023 & 0.85 & 0.88 & 0.83 \\
方法B & 2024 & 0.89 & 0.91 & 0.87 \\
\textbf{本文方法} & 2025 & \textbf{0.92} & \textbf{0.94} & \textbf{0.90} \\
\bottomrule
\end{tabular}
\end{table}

\subsection{消融实验}
分析各模块/超参数对性能的影响，如表~\ref{tab:ablation}。

\begin{table}[htbp]
\centering
\caption{消融实验结果}
\label{tab:ablation}
\begin{tabular}{lccc}
\toprule
配置 & Acc & AUC & F1 \\
\midrule
基线 & 0.85 & 0.88 & 0.83 \\
+ 模块A & 0.88 & 0.90 & 0.86 \\
+ 模块B & 0.90 & 0.92 & 0.88 \\
完整方法 & 0.92 & 0.94 & 0.90 \\
\bottomrule
\end{tabular}
\end{table}

% ========== 5. 结果与讨论 ==========
\section{结果与讨论（Results and Discussion）}
对实验结果进行深入分析，可包含：可视化结果、失败案例分析、方法优势与局限性讨论。

\begin{figure}[htbp]
\centering
\includegraphics[width=0.8\textwidth]{example-image-a}
\caption{实验结果可视化示例。请替换为实际图片。}
\label{fig:result}
\end{figure}

% ========== 6. 结论 ==========
\section{结论（Conclusion）}
总结本文的主要工作，强调创新点和实验结论。可简要提及未来研究方向。

% ========== 致谢（可选） ==========
\section*{致谢（Acknowledgments）}
感谢资助项目、设备支持、同行讨论等。

% ========== 参考文献 ==========
\bibliography{references}

% ========== 附录（可选） ==========
\appendix
\section{补充材料}
[请将补充实验、推导、伪代码等放在附录中。]

\end{document}
